% =========================================================================
% APA 7th Edition LaTeX Template for Academic Essays and Papers
% =========================================================================
% Document class options:
% - [stu]: Student paper format (default, no running head, includes course info)
% - [man]: Manuscript format (standard for journal submissions, includes running head & abstract)
% - [jou]: Journal format (two-column, compact layout)
% =========================================================================
\documentclass[stu, 12pt, american]{apa7}

% --- Essential Packages ---
\usepackage[utf8]{inputenc}
\usepackage{babel}
\usepackage{csquotes}

% --- APA 7th Edition Citations (BibLaTeX + Biber) ---
\usepackage[style=apa, backend=biber]{biblatex}
\addbibresource{references.bib}

% --- Document Metadata ---
\title{Your Paper Title: Capitalize the Main Words of Your Subtitle}
\shorttitle{APA 7 LaTeX Template} % Required if using [man] (manuscript) mode

\author{Your Name}
\affiliation{Department of Your Major, Your University Name}

% The following fields are specific to student [stu] mode:
\course{SUBJ 1234: Course Name}
\professor{Instructor Name}
\duedate{Month DD, YYYY}

% --- Abstract (Required for [man] mode, ignored in [stu] mode) ---
\abstract{
  This is a sample abstract for the document. An abstract is a brief, comprehensive summary of the contents of the paper. It is typically required for professional manuscripts (using the \texttt{[man]} document class option) but is usually omitted for student papers (using the \texttt{[stu]} option) unless requested by the instructor. The abstract should be a single paragraph, double-spaced, and between 150 to 250 words.
}

% =========================================================================
% Document Body
% =========================================================================
\begin{document}

\maketitle % Generates the cover page automatically based on document class settings

% -------------------------------------------------------------------------
% Main Content Section
% -------------------------------------------------------------------------
\section{Introduction}
Start writing the introduction of your essay here. The title of the paper is automatically centered and bolded at the top of the first page of text as a Level 1 heading, so you do not need to repeat it manually.

In LaTeX, you can cite sources dynamically using the \texttt{biblatex} package. For instance, you can cite a source parenthetically \parencite{sample_article} or narratively as described by \textcite{sample_book}.

APA style requires double-spacing, 1-inch margins, and page numbers in the top-right corner. The \texttt{apa7} class configures these margins and line spacing automatically.

\section{First Major Topic}
This is a Level 1 heading. Under APA 7 rules, headings are structured hierarchically:
\begin{itemize}
    \item \textbf{Level 1:} Centered, Bold, Title Case Heading.
    \item \textbf{Level 2:} Flush Left, Bold, Title Case Heading.
    \item \textbf{Level 3:} Flush Left, Bold Italic, Title Case Heading.
\end{itemize}

\subsection{Subtopic A}
This is a Level 2 heading. It is formatted left-aligned, bold, and in title case.

\subsubsection{Sub-subtopic 1}
This is a Level 3 heading. It is formatted left-aligned, bold italic, and in title case.

\section{Conclusion}
To finish your document, write your concluding thoughts here. Make sure all citations correspond to entries in your bibliography file.

% -------------------------------------------------------------------------
% References Page
% -------------------------------------------------------------------------
\clearpage
\printbibliography % Generates the "References" section automatically on a separate page

\end{document}
